\title{Large Language Models Can Automatically Engineer Features for Few-Shot Tabular Learning}

\begin{abstract}
Large Language Models (LLMs) have demonstrated extraordinary proficiency in addressing complex and novel reasoning tasks, suggesting significant promise for applications in tabular data analysis—an area fundamental to numerous practical scenarios. In this study, we introduce a new in-context learning paradigm called \textsf{FeatLLM}, which leverages LLMs as feature engineering agents to construct a transformed dataset that enhances tabular prediction performance. The newly created features serve as input to a straightforward downstream learning algorithm, such as linear regression, thereby supporting effective few-shot learning. Crucially, \textsf{FeatLLM} employs only this lightweight predictive model—augmented by the extracted features—during the inference stage. Unlike conventional LLM-based solutions, our framework circumvents the necessity to interact with the LLM for every data sample during inference. Furthermore, \textsf{FeatLLM} is operable using only API-level LLM access and is unaffected by prompt length constraints. Empirical evaluations across a variety of tabular datasets spanning multiple application domains show that \textsf{FeatLLM} formulates high-quality feature representations and rules, achieving average improvements of 10\% in comparison to established approaches like TabLLM and STUNT.
\end{abstract}

\section{Introduction}
Large Language Models (LLMs) have recently attained state-of-the-art performance in providing relevant outputs for previously unencountered tasks, all without relying on fine-tuning for specific tasks~\cite{creswell2022selection,imani2023mathprompter,taylor2022galactica}. LLMs, pretrained on vast amounts of textual data, encapsulate broad prior knowledge that can effectively substitute for human expertise across various domains~\cite{Genomes2021,singhal2023large,wilcox2003role}. This substitution has enabled the automation of complex processes, including tasks such as dialogue interpretation~\cite{finch2023leveraging} and automated code fixing~\cite{fan2023automated}. By integrating a task prompt with a handful of illustrative examples, these models are capable of comprehending task context and highlighting salient features and examples~\cite{brown2020language,zhao2023survey}, which reveals their aptitude for serving as automated feature engineers capable of analyzing and synthesizing data attributes.

Expanding on these competencies, LLMs’ background knowledge and reasoning power have recently found application within tabular data learning. For example, insights—such as age-related risk for specific diseases—are valuable for disease prediction problems. Emerging research exploits these capacities by constructing task and feature descriptions in natural language, transforming the original data accordingly, and processing this serialized representation using LLMs~\cite{dinh2022lift,wang2023anypredict}. This method is typically augmented either by in-context learning strategies~\cite{nam2023semi} or efficient tuning of model parameters~\cite{dinh2022lift,hegselmann2023tabllm}. Results to date indicate that incorporating the contextual understanding embedded in LLMs not only surpasses standard tabular learning methods but is also particularly beneficial in settings with limited labeled data~\cite{hegselmann2023tabllm}.

Prevailing LLM-based techniques for tabular data analysis face notable challenges. First, methods that make end-to-end predictions demand at least one LLM evaluation for each data point, substantially increasing computational load. High predictive accuracy is usually contingent on fine-tuning the LLM, but such adjustments are not practical for most LLMs lacking public training infrastructure or tunable APIs. This limitation is especially problematic as recent state-of-the-art LLMs generally restrict access to inference-only APIs~\cite{anil2023palm,openai2023gpt4}. Furthermore, these techniques often struggle with long input prompts~\cite{liu2023lost}: when tabular datasets have many features, the serialized text fed to the model can become unwieldy, thereby rendering the method impractical or degrading its effectiveness. Collectively, these issues impede the deployment of LLM-based tabular learning strategies in practical settings.

To address these shortcomings, our method reframes the use of LLMs from prediction engines to advanced feature engineers, unlocking their latent knowledge more efficiently. We propose a new in-context learning paradigm, \textsf{FeatLLM}, which focuses on identifying the “criteria” guiding LLM prediction logic. In scenarios such as disease prediction, the LLM can elucidate and formalize the decision rules that clarify how certain feature combinations correlate with outcomes. These criteria are translated into newly derived features that supplant the originals for subsequent learning tasks. With this setup, the heavy computational demands of sophisticated machine learning models are alleviated for inference purposes, thanks to the improved descriptive power of LLM-generated features—this streamlines and accelerates inference. Because LLMs can be steered via in-context learning (by providing example cases in their prompts), no model training is necessary to obtain effective features. Lastly, integrating strategies from ensemble learning~\cite{breiman2001random}, such as feature bagging, helps overcome difficulties posed by overly lengthy input prompts.

\textsf{FeatLLM} organizes the process of crafting new features from prompts into two core stages: (1) comprehending the task at hand by discerning the association between candidate features and the target variable; and (2) leveraging insights from this analysis, along with limited annotated instances, to formulate discriminative rules that separate the different categories. By guiding LLMs through this stepwise reasoning paradigm, the approach encourages the exclusion of non-informative features during rule formulation. The resultant rules are executed on the dataset (translating them as code), producing binary indicators that signal whether each instance complies with the respective rule. Subsequently, a simple classifier is fitted to predict class membership using these newly constructed features. To increase robustness, an ensemble methodology is adopted, wherein feature sets are regenerated several times using a feature bagging strategy. Selecting suitable feature and instance subsets for each ensemble iteration helps address the problem of oversized prompt inputs. As \textsf{FeatLLM} does not necessitate model training and maintains a low inference burden, it extends the practicality of LLM-driven feature engineering to a wide array of application domains.

Our experimental assessment of \textsf{FeatLLM} spans 13 tabular datasets in scenarios with scarce labeled data, evidencing both strong efficacy and consistency. The results demonstrate that LLMs can efficiently synthesize both prior domain knowledge and insights drawn from data, leading to the derivation of highly informative features. Compared to recent few-shot learning frameworks, our method achieves superior performance under various experimental conditions. The implementation is made publicly available through an anonymous GitHub repository at~\texttt{https://github.com/Sungwon-Han/FeatLLM}.

\section{Related Work}

\subsection{Few-Shot Learning with Tabular Data} The progression of few-shot learning has made it possible to derive robust and general representations from datasets, even when annotated examples are scarce~\cite{chen2018closer}. While initial research primarily centered on images, the principles of few-shot learning have also shown considerable adaptability to various other modalities, such as text, audio, and, most recently, tabular data~\cite{majumder2022few,nam2023stunt,schick2021exploiting}. The inherent limitation of having few labeled instances raises the danger of acquiring misleading or spurious patterns~\cite{bartlett2021deep}. To mitigate this, contemporary research has integrated supplementary informational resources or embedded domain knowledge, which introduce beneficial inductive biases into model learning. Some approaches achieve this by leveraging ample, easily accessible unlabeled datasets combined with strategic data augmentation to enhance semi-supervised learning~\cite{ucar2021subtab,yoon2020vime}, or by employing unsupervised, task-agnostic pre-training to enable strong initialization for downstream tasks~\cite{somepalli2021saint}. Methods under the unsupervised paradigm might involve altering or masking subsets of data followed by attempting to reconstruct the original information~\cite{arik2021tabnet,majmundar2022met} or adopting augmentation techniques for contrastive learning objectives~\cite{bahri2022scarf,chen2020simple}. Nevertheless, the diverse and mixed nature of tabular data creates obstacles for designing augmentation techniques that are effective across datasets, and these techniques often fall short in scenarios where labeled samples are especially limited~\cite{nam2023stunt}.

A line of recent investigation has involved training models on an extensive collection of disparate tabular datasets—extending well beyond those immediately relevant to the downstream task—and subsequently applying transfer learning to new tasks~\cite{zhang2023generative,levin2022transfer,zhu2023xtab}. As an illustration, TransTab~\cite{wang2022transtab} utilizes transformer-based architectures to understand semantic associations between columns using column names, rather than just focusing on the tabular values. Such representations capture transferable knowledge from diverse columns and tables, making zero-shot or few-shot inference on previously unseen tabular tasks feasible. Similarly, TabPFN~\cite{hollmann2022tabpfn} synthesizes a variety of data distributions to pre-train a Bayesian neural network model. Upon completion of pre-training, inference is achieved by simultaneously providing the model with training and testing examples from the new task, bypassing any further fine-tuning. The prior information distilled from the heterogeneous datasets enables this method to not only outperform standard tree-based models, but also to be especially effective in few-shot learning contexts.

\subsection{Language-Interfaced Tabular Learning} Integrating large language models (LLMs)~\cite{anil2023palm,openai2023gpt4} represents another significant strategy for incorporating prior knowledge into tabular learning. With their extensive exposure to varied textual data, LLMs have proven capable of adapting to new, unseen tasks and are showing effectiveness in an array of fields~\cite{brown2020language}. Recently, research has investigated the potential of utilizing the comprehensive background knowledge of LLMs to enhance tabular data processing. This is typically accomplished by converting tables into text format and embedding them within the input prompts for LLMs~\cite{dinh2022lift,hegselmann2023tabllm,wang2023anypredict}. Such prompts generally combine the tabular entries, feature explanations, and task details, enabling models to deduce relations among tasks and features to carry out inference. Methodological advances in this area include the fine-tuning of models using parameter-efficient techniques like LoRA~\cite{hu2021lora} or IA3~\cite{liu2022few}, as well as the application of in-context learning by providing the model with a handful of examples in the prompts, obviating the need for retraining~\cite{dinh2022lift,hegselmann2023tabllm,nam2023semi,slack2023tablet}.

Although prior techniques have demonstrated success in advancing few-shot tabular learning, their reliance on routing every inference through the LLM introduces constraints for real-world industrial deployment. This work seeks to move away from the traditional black-box paradigm where predictions for individual samples are handled exclusively by an LLM in a sequential manner. Instead, the objective is to harness the LLM’s capacity for directly identifying the distinguishing criteria or rule sets that pertain to each class. With \textsf{FeatLLM}, these rules are extracted and systematically translated into program code, yielding newly engineered features. Such an approach enables subsequent predictions to bypass the LLM entirely, while in-context learning is utilized to harvest relevant rules using both prior knowledge and the information derived from the limited available examples.

\section{Methods}
\textsf{FeatLLM} is structured to systematically mine the “rules”—that is, the defining conditions for each target class—based on the limited examples provided, as visualized in Figure~\ref{fig:main_model}. The system leverages the LLM to generate these rule sets, which are then translated into logical program code. From here, new binary indicators are created for each data sample, encoding whether each particular rule holds true for that instance. These binary indicators are then subjected to a dot product with learnable, non-negative coefficients to generate the class probability estimates. During training, the model learns which engineered features are most salient for accurate class discrimination (see Section~\ref{Sec:training}). This rule extraction and feature generation process is repeated with the use of bagging, and the resulting outputs are integrated via ensemble methods. A more detailed discussion of problem setting and the specifics of each processing stage follows.

\vspace*{-0.12in}
\paragraph{Problem formulation.} Consider a tabular dataset containing $N$ labeled instances, denoted as $\mathcal{D} = \{(\mathbf{x}^i, \mathbf{y}^i)\}_{i=1}^{N}$. Each entry in $\mathcal{D}$ has $d$ features, meaning each $\mathbf{x}^i$ corresponds to a vector in $\mathbb{R}^d$, and every feature has an associated natural language name, summarized as $F = \{f_j\}_{j=1}^{d}$. Feature definitions are supplied separately. In a classification context, each label $\mathbf{y}^i$ belongs to a designated set of class labels. For $k$-shot learning scenarios, we select $k$ ($<N$) labeled data points at random to train the model.

\subsection{Prompt Design for Extracting Rules}
\label{Sec:prompt_design}
To prompt the LLM to derive rules following a more robust and interpretable reasoning path, we structure the interaction to reflect an expert human’s approach to analyzing tabular data. The prompt comprises three key sections, as illustrated in Fig.~\ref{fig:prompt_for_rule} and further elaborated in Appendix~\ref{sec:example_rule}:

\vspace*{-0.12in}
\paragraph{Basic information description.} This section communicates the essential context needed to address the learning task (refer to the orange text in Fig.~\ref{fig:prompt_for_rule}). It encompasses a statement of the task including the set of possible labels, detailed descriptions of each feature, and select illustrative examples. Task formulations are cast in the form of explicit questions, such as ``Does this patient's myocardial infarction complications data indicate chronic heart failure? Yes or no?" (see Appendix~\ref{sec:task_desc} for further instances). Feature descriptions provide information about type and, when available, their specific definitions; categorical variables may include example categories. For demonstration, a few labeled samples from the training set are presented as serialized text together with their actual labels. Serialization adopts the protocol established in earlier works~\cite{dinh2022lift,hegselmann2023tabllm}:

\begin{align}
&\text{Serialize}(\mathbf{x}^i,\mathbf{y}^i, F) = \nonumber \\ 
&``\ f_1\ \text{is}\ \mathbf{x}^i_1.\ ... \ f_d\  \text{is}\  \mathbf{x}^i_d.\ \ \text{Answer:}\  \mathbf{y}^i\ ”
\end{align}

\vspace*{-0.12in}
\paragraph{Reasoning instruction.} Our goal is to improve the reasoning capabilities of LLMs by supplying them with tailored reasoning instructions (see the blue text in Fig.~\ref{fig:prompt_for_rule}). Each instruction begins with an opening statement reminiscent of the chain-of-thought framework~\cite{wei2022chain}. We then structure the reasoning procedure for the LLM into two distinct stages. First, the model is prompted to determine the causal links or patterns between various features and the stated task, while deliberately omitting consideration of any sample demonstrations; instead, this step capitalizes on the model’s world knowledge or innate reasoning skills, ensuring comprehensive leverage of prior understanding relevant to the problem. Subsequently, in the second phase, the LLM synthesizes information gleaned from both the initial analysis and provided demonstrations to generate classification rules specific to each class. Segmenting the reasoning process in this fashion helps the model avoid picking up misleading associations from irrelevant data columns, thus promoting attention toward more pertinent attributes. To strike a balance and avoid producing too few rules (which may cause underfitting) or an unwieldy prompt with too many, we standardize the process by capping the number of rules generated per class at 10\footnote{Analysis on the number of rules can be found in Section~\ref{sec:analysis}.}. Additionally, during rule formulation, the LLM is directed to consult feature type information included in the basic descriptive section to avert generating rules that are incompatible with feature types.

\vspace*{-0.12in}
\paragraph{Response instruction.} To ensure the inferred rules can be easily interpreted and utilized, we instruct the LLM on the precise structure to follow when composing its responses (see the yellow text in Fig.~\ref{fig:prompt_for_rule}). The prompt details explicit formatting conventions for responses specific to each answer class (i.e., how to organize the response), as well as the required format for individual rules (i.e., template for rule expression). These response instructions permit the LLM to choose formatting that aligns with the characteristics of each feature. In the absence of further constraints, the LLM is also able to logically combine multiple rules via operators such as AND or OR. An illustrative response can be found in Appendix~\ref{sec:outcome-rule}.

\subsection{Inferring Class Likelihood via Rules}
\label{Sec:training}

\paragraph{Parsing rules for feature generation.} In this stage, we transform the previously generated rules into new binary features. For each class, a binary indicator is produced, signaling whether a given sample complies with the specified rules for that class (coded as '0' or '1'). These indicators are incorporated into the prediction of a sample's class membership. Nonetheless, because these rules are derived from natural language via an LLM, a dedicated parsing strategy is necessary to enable the automatic construction of such features. Although we impose strict response and rule formatting requirements during rule creation for greater parsing consistency, outputs from the LLM may still contain inconsistent syntactic modifications~\cite{kovavc2023large}. For example, an LLM might replace symbols such as ``$>$” or ``$<$” with their word equivalents, like ``greater than" or ``less than," which introduces additional parsing complexity.

To manage the intricacies of extracting structured data from inconsistent or noisy textual rules, we avoid intricate custom programmatic parsing and instead employ the LLM itself. The LLM's robust semantic comprehension—despite potential syntactic variations—as well as its proficiency in generating executable code (benefiting from extensive training on code), are both utilized in our approach. We format a prompt by supplying the function name, detailed input/output descriptions, and the corresponding rules, which is then submitted to the LLM. Illustrative prompt samples and the results generated are provided in Figures~\ref{fig:prompt_for_parsing} and ~\ref{sec:example_parsing}-\ref{sec:example_parse_outcome} in the Appendix. The resulting code is implemented through Python's \texttt{exec()} method to convert the data accordingly.

\vspace*{-0.12in}
\paragraph{Inferring class likelihood.} When binary features are generated from rules associated with each class, a basic technique to estimate the likelihood of a sample belonging to a specific class involves tallying the number of rules from each class that the sample fulfills (i.e., summing the corresponding binary features for each class). Nevertheless, the influence of different rules and features varies, making it crucial to determine their relative importance based on the training set. We propose to accomplish this by fitting a traditional machine learning (ML) model to the binary features linked to every class. For instance, a bias-free linear model can be employed. Consider a sample $\mathbf{x}^i$, for which the class-$k$ rule-derived binary feature is $\mathbf{z}_k^i \in \{0, 1\}^R$, where $R$ is the quantity of rules available per class and $c$ denotes the number of classes. The class probability vector $\mathbf{p}^i$ is calculated using:

\begin{align}
\text{logit}_k^i &= \max(\mathbf{w}_k, 0) \cdot \mathbf{z}_k^i, \nonumber \\
\mathbf{p}^i &= \text{Softmax}({[\text{logit}_{1}^i, ..., \text{logit}_{c}^i]}). \label{eq:infer_prob}
\end{align}

Here, $\mathbf{w}_k$ refers to the set of trainable weights in the linear model, indicating the significance of each rule-derived feature. Restricting these weights to non-negative values guarantees that none of the features extracted by the LLM will diminish the probability of a given class during learning. The resulting logit values are normalized into class probabilities via the softmax operation. The parameters $\mathbf{w}_k$ are optimized through cross-entropy loss, utilizing few-shot training data. 

\vspace*{-0.12in}
\paragraph{Ensembling with bagging.} To improve robustness, we repeat the full pipeline multiple times and aggregate predictions across an ensemble of inference models. For $T$ repetitions, we obtain a set of trained weights $\{\mathbf{w}[t]\}_{t=1}^T$; then, the prediction for each sample is made by averaging the class probability vectors $\mathbf{p}^i$, each computed using a different $\mathbf{w}[t]$ as in Eq.~\ref{eq:infer_prob}.

To incorporate stochasticity into rule generation for ensemble learning, we set a relatively high temperature during LLM inference. Additionally, we shuffle the sequence of few-shot examples, as prior studies have found that their ordering can influence the LLM's response~\cite{lu2022fantastically}\footnote{Appendix~\ref{sec:diversity} provides an analysis of rule diversity.}. To further harness prediction variability for improved model resilience, we employ bagging, randomly selecting subsets of features or instances in each run. This is particularly advantageous when training datasets are large or contain numerous features. The ensemble strategy proposed here confers two principal benefits. First, even if the LLM yields suboptimal rules in certain runs, correct rules generated elsewhere can counterbalance these errors, thereby strengthening the system's overall robustness. This property aligns with the concept of LLM self-consistency~\cite{wang2022self}, where rules that repeatedly appear across trials have a higher likelihood of accuracy. Second, the ensemble framework mitigates LLM prompt size limitations: for tabular datasets exceeding input capacity, bagging reduces input dimensionality and safeguards performance. While any single generated rule set may fail to encapsulate all feature dependencies, aggregating diverse weak learners across multiple runs diminishes the adverse effects of partial feature or instance coverage in individual trials.

\section{Experiments}
We assess \textsf{FeatLLM} on a range of tabular datasets and provide several analyses to elucidate various aspects of our framework. More extensive experimental results—including those involving different LLMs, in-depth qualitative studies, and further analyses—can be found in the Appendix due to space limitations.

\subsection{Performance Evaluation}
\paragraph{Datasets.}
Our study leverages 13 datasets designed for binary or multi-class classification: (1) Adult~\cite{asuncion2007uci}, which predicts if a person’s income exceeds \$50,000 per year; (2) Bank~\cite{moro2014data}, which identifies customers likely to subscribe to a term deposit; (3) Blood~\cite{yeh2009knowledge}, aimed at predicting whether donors will make additional donations; (4) Car~\cite{kadra2021well}, which classifies automobile quality; (5) Communities~\cite{misc_communities_and_crime_183}, used for estimating regional crime rates; (6) Credit-g~\cite{kadra2021well}, focused on assessing whether someone represents a good or bad credit risk; (7) Diabetes\footnote{\texttt{kaggle.com/datasets/uciml/pima-indians-diabetes-database}}, which determines the likelihood of diabetes in individuals; (8) Heart\footnote{\texttt{kaggle.com/datasets/fedesoriano/heart-failure-prediction}}, which predicts the presence of coronary artery disease; and (9) Myocardial~\cite{misc_myocardial_infarction_complications_579}, which forecasts chronic heart failure in patients.

Additionally, we incorporate newer datasets and synthetic datasets that are less commonly featured in literature and absent from the pre-training of LLMs: (10) Cultivars, used for predicting the expected grain yield of specific soybean cultivars\footnote{\texttt{archive.ics.uci.edu/dataset/913}}; (11) NHANES, for classifying an individual's age group based on provided records\footnote{\texttt{archive.ics.uci.edu/dataset/887}}; (12) Sequence-type, aimed at determining whether a given sequence is arithmetic, geometric, Fibonacci, or Collatz; (13) Solution-mix, which assesses if the concentration percentage from a blend of four solutions surpasses 0.5. The first two datasets were only recently made available through the UCI Machine Learning Repository after September 2023, guaranteeing their exclusion from the training data of the LLM API (gpt-3.5-turbo-0613) utilized in our study. The last two datasets are artificially generated, preventing any possibility of memorization by the LLM API and thereby maintaining fair evaluation.

There is considerable diversity in both dataset scale and level of complexity, as outlined in Table~\ref{Tab:basic-desc}. Each dataset offers explicit naming and explanations for every attribute. Some datasets, such as `Communities' and `Myocardial,' encompass more than 100 features, presenting significant obstacles for LLMs constrained by limited context windows.

\vspace*{-0.12in}
\paragraph{Baselines.}
Our evaluations involve comparison with nine baseline methods. The first three baselines reflect standard models in tabular data analysis: (1) Logistic regression (LogReg), (2) XGBoost~\cite{chen2016xgboost}, and (3) RandomForest~\cite{ho1995random}. The fourth and fifth baselines are tailored to scenarios utilizing large unlabeled datasets, including (4) SCARF~\cite{bahri2022scarf} and (5) STUNT~\cite{nam2023stunt}. It should be recognized that obtaining abundant unlabeled data is often impractical in many applied settings. A contemporary alternative is (6) TabPFN~\cite{hollmann2022tabpfn}, which leverages pretraining on synthetic datasets with varying distributions. The remaining three baselines utilize large language models: (7) In-context learning (In-context)~\cite{wei2022emergent}, (8) TABLET~\cite{slack2023tablet}, and (9) TabLLM~\cite{hegselmann2023tabllm}. In-context learning operates by embedding a limited number of training examples into the input prompts without additional model tuning. TABLET improves upon this by supplying auxiliary guidance via external classifiers—such as rule-sets and prototypes—within the prompts to strengthen inference performance. TabLLM adapts the LLM via parameter-efficient tuning methods, such as IA3~\cite{liu2022few}, based on a small sample of training instances.

\vspace*{-0.12in}
\paragraph{Implementation details.}
\textsf{FeatLLM} is built upon GPT-3.5 as its foundational LLM, but it is purposefully structured to support various LLM choices (see Appendix~\ref{sec:backbones} for comparative analyses across different backbones). The LLM’s inference procedure employs a temperature setting of 0.5 and retains the API’s default top-p parameter of 1. For rule extraction, we designate 20 ensembles and 10 rules. The influence of hyper-parameter choices can be found in Figure~\ref{fig:hyperparameter2} and is further explored in Appendix~\ref{sec:hyper-parameter}. The linear model is optimized using Adam with a 0.01 learning rate over 200 epochs. Selection of the optimal training epoch relies on $k$-fold cross-validation. For baseline re-implementations, GPT-3.5 is used for in-context learning, while TabLLM is paired with T0~\cite{sanh2022multitask} as prescribed in its original methodology. Importantly, TabLLM constrains LLM options to those with open, public checkpoints, making GPT-3.5 inaccessible within this framework. For conventional machine learning approaches such as LogReg, XGBoost, and RandomForest, parameter tuning utilizes Grid-search in conjunction with $k$-fold cross-validation. We select $k$ as either 2 or 4, ensuring each class is represented in the training data. Comprehensive implementation details are included in Appendix~\ref{sec:baselines}.

\vspace*{-0.12in}
\paragraph{Results.}
Tables~\ref{tab:main1} and \ref{tab:main2} display the comparative results across multiple datasets. Each experiment is conducted three times, with varying random seeds for splitting training and testing data; we report both the mean and standard deviation of the resulting AUC metrics. When compared to other baseline methods, \textsf{FeatLLM} either attains the best results or occupies the second highest rank consistently. With only 4 training samples per class, our approach is able to deliver performance on par with traditional tabular models trained with 32 samples (refer to Fig.~\ref{fig:compares}). Although STUNT and TabPFN achieved noticeable gains on selected datasets, their improvements over standard machine learning baselines remain modest overall. In addition, LLM-based models—including in-context learning, TABLET, and TabLLM—struggled to process tabular data characterized by high-dimensional feature spaces. By integrating bagging and ensemble methodologies, our framework retains strong efficacy even on datasets with a large number of features. It is also important to highlight that the bagging and ensemble paradigm proposed here could theoretically be incorporated into other LLM-driven frameworks, but this would result in doubling the inference passes per data instance, dramatically increasing computational cost and making it less feasible in practice.

\subsection{Analysis \& Discussion}
\label{sec:analysis}
\paragraph{Ablation study.} To evaluate the contribution of individual framework components, we conduct a series of ablation experiments that sequentially remove key features: (1) \textsf{-Tuning}: the linear model does not undergo weight tuning, using a simple aggregation of binary features to produce the logit; (2) \textsf{-Ensemble}: eliminates the ensemble phase; (3) \textsf{-Description}: omits descriptive text for features; (4) \textsf{-Reasoning}: skips the Step-1 phase involved in producing reasoning instructions for rule creation.

Results from these ablation scenarios, as provided in Table~\ref{tab:ablation}, indicate the average change in AUC across datasets. The removal or modification of any major component invariably reduces performance. Of note, the impact of weight tuning becomes more pronounced with an increased number of training examples, which implies that tuning has greater value as more data allows better assessment of rule significance. Conversely, when shot counts are low, both feature descriptions and reasoning instructions contribute more heavily to model effectiveness, suggesting that leveraging LLM priors is critical under data-limited conditions. Consistently, our ensemble strategy delivers robust gains irrespective of the experimental setting.

\vspace*{-0.12in}
\paragraph{Training \& inference time comparison.} We evaluate and contrast the duration required for training and inference across various methods. All experiments utilize the Adult dataset, which consists of a training set with 16 samples. One A100 GPU is employed for all approaches, apart from TabLLM, which requires four A100 GPUs to facilitate model parallelism. The detailed runtime outcomes are shown in Table~\ref{tab:cost}. It is noteworthy that \textsf{FeatLLM} achieves low inference latency, matching that of standard approaches. In contrast, inference times for other LLM-powered techniques, including In-context learning, TABLET, and TabLLM, are substantially higher. Regarding training time, \textsf{FeatLLM} is on par with the other LLM-based strategies; it necessitates merely 30 API calls, keeping computational costs manageable. Additionally, these API requests can be dispatched in parallel, minimizing bottlenecks.

\vspace*{-0.12in}
\paragraph{Quality of features generated by \textsf{FeatLLM}.}
We perform an exploratory experiment to assess how informative the rule-driven features produced by \textsf{FeatLLM} are for real-world problem solving. For this purpose, we compute the mean AUROC between each new feature and the corresponding ground-truth labels and determine the proportion of features whose AUROC surpasses 0.5 for each dataset. The results, presented in Table~\ref{tab:quality}, demonstrate that most generated rules have substantial relevance to the prediction target, thereby contributing significant information toward solving the task (see the ``Before tuning'' column). Furthermore, when optimizing the linear model, rules offering minimal value (identified as having learned weights below a set threshold) are discarded automatically (as indicated in the ``After tuning'' column). In this context, the threshold value is fixed at 0.1, based on the analysis of the weight distribution.

\vspace*{-0.12in}
\paragraph{Effect of adding spurious correlations.} To evaluate the capacity of different methods to disregard irrelevant spurious correlations, particularly under limited labeled data scenarios, we carried out a targeted analysis. Specifically, we selected the Heart dataset and augmented it by randomly introducing columns from the Adult dataset that do not pertain to the Heart task. We then monitored how model performances shifted as we increased the number of these extraneous columns. To ensure consistency, each approach was trained with exactly 8 labeled examples, omitting any unlabeled data; as a result, approaches such as SCARF and STUNT were not included in this comparison.

In Figure~\ref{fig:spurious}, we display the changes in accuracy attributable to the introduction of spurious correlations. Among the evaluated methods, \textsf{FeatLLM} exhibited the smallest decline in performance as the number of unrelated columns grew. This resilience can be attributed to the design of our approach, which explicitly encourages rules that connect features to the main task in a manner grounded in prior knowledge, thereby discouraging rules built from irrelevant features and strengthening robustness. Quantitatively, we found that just 1-2\% of the induced rules incorporated the spurious columns. Nevertheless, we also note that if random columns from data closely related to the main task are added, \textsf{FeatLLM} can mistakenly internalize spurious patterns and misconstrue the causal structure between task and input features (refer to Appendix~\ref{sec:spurious-appendix} for details).

\vspace*{-0.12in}
\paragraph{Hyper-parameter analysis}
\textsf{FeatLLM} operates with two primary hyper-parameters: the ensemble count, and the number of rules generated per class in each LLM interaction. To understand their effect on the model, we systematically varied these settings in experiments. Findings indicate that while increasing the number of ensembles generally enhances performance, there is a diminishing return when the count approaches 20, beyond which improvements plateau (refer to Fig.~\ref{fig:appendix-hyperparameter1} in the Appendix). Concerning the rule count per class, no notable statistical difference was found across values; nonetheless, setting this number to 10 offers a balanced compromise between concise prompts and robust predictive accuracy (see Fig.~\ref{fig:hyperparameter2}). We hypothesize that selecting a higher number of rules adds to the input’s dimensionality, thus augmenting information for the model. However, in low-shot scenarios, this also potentially leads to greater overfitting.

\section{Conclusion}
We introduce \textsf{FeatLLM}, a novel approach that elevates the standards for LLM-driven few-shot learning on tabular datasets. Distinctively, \textsf{FeatLLM} leverages LLMs not merely for predicting labels, but for formulating underlying decision criteria akin to feature engineering practices. By integrating automated rule extraction with a bagging-based ensemble mechanism, the method ensures low inference overhead, is fully operational using API interfaces without model retraining, and surpasses the limitations imposed by feature dimensionality. Our extensive evaluations demonstrate that \textsf{FeatLLM} provides strong predictive capabilities, representing a robust and efficient alternative for real-world tabular tasks with minimal data availability.

Although \textsf{FeatLLM} is specifically tailored for low-shot learning scenarios, future work will target expanding its feature creation mechanisms to accommodate larger datasets, explore more diverse feature representations beyond rule-based ones, and enhance its interpretability, thereby broadening its applicability to a greater variety of domains.

\section{Broader Impact} The development of \textsf{FeatLLM} seeks to harness LLMs’ inherent knowledge to address tabular learning challenges that surpass conventional supervised approaches, particularly under limited data conditions. By utilizing this external knowledge, the framework provides meaningful advancements in reducing the risk of overfitting with sparse labels. This approach is poised to be highly advantageous for domain practitioners in sectors such as finance and healthcare—fields where acquiring labeled data can be prohibitively resource-intensive and LLMs’ pretrained expertise (e.g., from finance or healthcare texts) is directly relevant. Moving forward, it is crucial to systematically assess the influence of biases and misinformation present in the LLMs’ foundational knowledge, since these factors could substantively shape—or even outweigh—predictions relative to available ground truth data. Such analyses are vital for ensuring safe and responsible adoption of the method at scale.

\end{document}